% Запасные слайды презентации. Содержимое при необходимости
% адаптируется под конкретную защиту.

\begin{frame}
  \frametitle{Дополнительный материал:~замечания рецензентов}
  \begin{itemize}
    \item Замечание -- ответ
    \item Замечание -- ответ
    \item Замечание -- ответ
    \item Замечание -- ответ
    \item Замечание -- ответ
  \end{itemize}
\end{frame}

\begin{frame}
  \frametitle{Научная новизна}
  В~работе впервые:
  \begin{itemize}
    \item предложена математическая модель представления многопоточной
      программы в~виде ЧУМ операций с~памятью, сводящая задачу анализа
      консистентности к~построению линейных порядков, с~разделением
      ограничений порядка $P$ и~зависимости операций $D$;
    \item разработаны алгоритмы эффективного построения всех линейных
      порядков ЧУМ с~учётом классов эквивалентности перестановок, которые
      порождают все линейные порядки \\textbf{с~учётом аксиом модели
      памяти} и~не~строят заведомо недопустимых кандидатов;
    \item разработана структура данных --- граф состояний (State~Graph)
      --- для компактного и~полного представления результатов исполнения;
    \item предложены алгоритмы верификации на~основе линейных порядков
      и~на~основе графа состояний.
  \end{itemize}
\end{frame}
\note{
  Проговаривается вслух научная новизна.
}

\begin{frame}
  \frametitle{Научная и~практическая значимость}
  \begin{description}
    \item[Теоретическая значимость] --- математическое описание
      многопоточных программ с~учётом моделей памяти, алгоритмы
      построения линейных порядков, структура данных графа состояний;
    \item[Практическая значимость] --- инструментарий верификации
      подсистемы памяти, снимающий ограничения существующих подходов
      по~типам операций и~пересечению адресов и~применимый к~программам
      с~большим числом конкурирующих обращений.
  \end{description}
\end{frame}
\note{
  Проговариваются вслух научная и~практическая значимость.
}

\begin{frame}
  \frametitle{Резервный слайд}
  Здесь может быть размещён дополнительный иллюстративный материал
  (граф состояний, литмус-тесты, результаты экспериментов).
\end{frame}
